\documentclass[a4paper,14pt]{article}
\usepackage[utf8]{inputenc}
\usepackage[german,english,russian]{babel}
\usepackage{amssymb,amsmath,amsfonts,amsthm}
\usepackage{cmap}
\usepackage{fancyhdr}
\usepackage{float}
\usepackage{cite}
\usepackage{hyperref}
\usepackage[dvips,usenames]{color}
\usepackage{xcolor}
\usepackage[dvips]{graphicx}
\usepackage{tikz}
\usepackage{setspace}
\usepackage{enumerate}
\usepackage{tabularx}
\usepackage{colortbl}
\usepackage{multicol}
\usepackage{multirow}
\usepackage{chngpage}
\usepackage[mathscr]{eucal}
\usepackage{bbm}
\usepackage{mathbbol}
\usepackage{ifthen}
\bibliographystyle{article}

\usepackage[top=2cm,bottom=1cm,bindingoffset=0cm]{geometry}
\textheight=240mm
\textwidth=160mm
\oddsidemargin=0.1in
\evensidemargin=0.1in

%--------------------------------------------------------------------------------------------
\begin{document}
\sloppy
\renewcommand{\refname}{\small{References}}
\thispagestyle{fancy}
\fancyhead{}
\fancyhead[L]{\footnotesize{PhD Summer School in Discrete Mathematics}}
\fancyhead[R]{\footnotesize{Abstract}}
\fancyfoot{}
\renewcommand{\footrulewidth}{0.1 mm}

\begin{center}
  \textbf{$\alpha$-labeling of banana tree graphs}\\
  \vspace{\baselineskip}
  Borna Gojšić\\
  \emph{Faculty of Electrical Engineering and Computing, Zagreb, Croatia}\\
  borna.gojsic@fer.hr\\
  \vspace{\baselineskip}
  This is joint work with Anamari Nakić\\
\end{center}
\vspace{\baselineskip}

A banana tree graph, denoted by $B(n, k)$, is a tree constructed by connecting $n$ star graphs $S_k$ to a root vertex. A graceful labeling assigns the numbers $0$ through $|E|$ to vertices so that edge differences yield all values from $1$ to $|E|$ exactly once; an $\alpha$-labeling is a graceful labeling with an additional ordering property, motivated by work of Rosa~\cite{1}, and is central in graph decomposition theory. This article builds upon previous research on graceful labelings of banana trees and explores their $\alpha$-labelability. We utilize the symmetries of the $B(n,k)$ trees to give the number of all of the non-isomorphic $\alpha$-labelings based on specific subsets of them which we call {\it proper} and {\it central} $\alpha$-labelings. We classify all the $\alpha$-labelings for all the $B(n, 1)$, $B(1, k)$ and $B(2,k)$ trees. We present $\alpha$-labelings of the $B(n,k)$ tree, for $k \geq \frac{n}{2} + 1$ and $n \in \mathbb{N}$. Finally, we prove that the $B(n, k)$ tree is not $\alpha$-labelable for $1 < k \leq \frac{n}{2}$ and any $n \ge 3$, classifying all the $B(n, k)$ trees based on their $\alpha$-labelability.

\begin{thebibliography}{99}
  \bibitem{1} \small{A.~Rosa, On certain valuations of the vertices of a graph, in: Theory of graphs (Internat. Sympos., Rome, 1966), Gordon \& Breach, New York, 1967, 349--355.}
\end{thebibliography}

\end{document}
